\subsection{The actual polarized abelian quotient and its \(O\)-action}
\label{subsec:faithful-actual-abelian-quotient}

We retain the source field hypotheses: \(E\) is a CM number
field with host \(H\), \(O=\mathcal O_E\) has its specified
involution \(\bar a\), \(k\) is algebraically closed of
characteristic a prime \(p\), and \(t\) is a source
\(\mathrm{SAB}\) datum over \(X\), satisfying
\(\mathrm{PowFns}(t)\).
Let \(x:\operatorname{Spec}k\to X\) satisfy
\(\mathrm{GeomPF}(t,x)\).
Its witness is an actual torus extension
\[
       0\longrightarrow D(C)\stackrel i\longrightarrow G_x
                       \stackrel q\longrightarrow B\longrightarrow0,
 \tag{AQ1}\label{eq:faithful-geometric-quotient-extension}
\]
where \(C\) is finite free over \(\mathbb Z\), \(B/k\) is
abelian, \(q\) is a flat surjective quasi-compact group map,
and \(i\) represents its kernel.  The witness also gives a
line \(M\) on \(B\), its \(\mathrm{KTriv}\) property, and an
actual line isomorphism \(L_x\simeq q^*M\).
There is no \(O\)-action on \(B\) in this witness yet.

\begin{srlemma}[The normalized torus-quotient seesaw]
\label{lem:faithful-torus-quotient-seesaw}
Let \(Y/k\) be smooth and geometrically connected with a
\(k\)-point \(y_0\).  For the extension
\eqref{eq:faithful-geometric-quotient-extension}, let \(N\) be
a line on \(B\times Y\).
If \((q\times1)^*N\) is trivial on \(G_x\times Y\) and
\(N|_{B\times\{y_0\}}\) is trivial, then \(N\) is trivial.
\end{srlemma}
\begin{proof}
Choose a generating section \(u\) of \((q\times1)^*N\).
Translation in the \(D(C)\)-variable compares two sections
of the same pulled-back line; their ratio is a unit
\[
             r(\tau,g,y)=u(i(\tau)g,y)/u(g,y).
\]
The source's weight calculation shows that
\(r(\tau,g,y)=\tau(\chi)\) for one fixed \(\chi\in C\).
Here is that calculation on charts.
The base \(G_x\times Y\) is smooth and geometrically connected,
hence reduced and connected.
On an affine open of this base the torus coordinate ring is
a Laurent polynomial ring.  An invertible Laurent polynomial
over a reduced ring is, on each connected component, a unit
times a single monomial: multiplying it by its inverse and
comparing extremal terms forces this statement over every
prime quotient, and the resulting exponents form open-and-closed
pieces.  Thus the exponent is locally constant.
These local exponents agree on overlaps, and connectedness
makes it one fixed \(\chi\).
The coefficient is \(1\), by evaluating at \(\tau=1\).
This is the source's \srcid{exists\_weight} calculation; it
does not assume that the original section \(u\) was invariant.

Choose a generator \(s\) of \(N|_{B\times\{y_0\}}\).
The ratio of \(u(g,y_0)\) to \(q^*s(g)\) is a global unit
\(f(g)\) on \(G_x\), and
\[
                   f(i(\tau)g)=\tau(\chi)f(g).
\]
Consequently \(u(g,y)/f(g)\) is invariant under the torus.
If two points of \(G_x\times Y\) have the same image under
\(q\times1\), their \(G_x\)-components differ by a unique
kernel point after the kernel factorization.
The corrected generating section therefore has equal
pullbacks on the fpqc relation.
Faithfully flat descent of this section and its inverse
gives a generator of \(N\).
These are \srcid{exists\_unit\_weight} and
\srcid{exists\_gen\_of\_weight}.
\end{proof}

\begin{srproposition}[The actual \(\mathrm{AbQuot}\) datum]
\label{prop:faithful-actual-abelian-quotient}
From the displayed hypotheses one constructs:
\begin{enumerate}
\item a descended \(O\)-action on \(B\), with
\(q[a]_{G_x}=[a]_Bq\) for every \(a\in O\);
\item a rigidified, symmetric, cubical, primitive line \(P\)
on \(B\), adjoint for this action and its specified involution;
\item \(\mathrm{KTriv}(B,P)\) and an actual isomorphism
\(L_x\simeq q^*P\).
\end{enumerate}
Together with the original torus extension, these are exactly
the source's \(\mathrm{AbQuot}(t,x)\) fields.
\end{srproposition}
\begin{proof}
We first descend the action, then the adjoint line comparison,
and finally normalize the line.  These are the three separate
constructions in \srcid{AbelianQuotientOfPowerFns}.

For each \(a\in O\), the composite
\(D(C)\to G_x\stackrel{[a]}\longrightarrow G_x\to B\)
is zero by Lemma~\ref{lem:faithful-torus-proper-constant}.
The map \([a]q:G_x\to B\) is therefore constant on \(q\)-fibres.
Explicitly, if \(g_1q=g_2q\), then
\(g_1g_2^{-1}\) factors through \(i\); its composite with
\([a]q\) is zero.  Hence the two values agree.
The effective-epimorphism property of the flat surjective
quasi-compact \(q\) gives a unique morphism \([a]_B\) with
\[
                         q[a]_B=[a]_{G_x}q
\]
in the source's left-to-right composition notation.
This is \srcid{exists\_desc}.
It preserves the identity; rigidity for morphisms between
abelian varieties makes it a group homomorphism.
The equations for \(1\), products, and sums follow by
precomposing with the epimorphism \(q\) and using the equations
on \(G_x\).  Thus it is the actual descended \(O\)-action
\srcid{exists\_act}.

For the next step write \(m=[M]\) in the Picard group, and
write \(\lambda_m(u,v)\) for its normalized second difference:
\[
 [M_{u+v}]-[M_u]-[M_v]+[M_0].
\]
Set, on any test scheme,
\[
 D_a(u,v)=\lambda_m([a]_B u,v)
                   -\lambda_m(u,[\bar a]_Bv).
\]
It commutes with pullback, vanishes on either zero axis, and
\[
                         D_a(qg,qh)=0,
 \tag{AQ2}\label{eq:faithful-quotient-adjoint-pull}
\]
because \(q^*M\simeq L_x\), the actions commute with \(q\), and
\(L_x\) has the actual adjoint isomorphism of the original
\(\mathrm{SAB}\) datum.

Pull this class first to \(B\times G_x\).
Its pullback along \(q\times1\) is zero by
\eqref{eq:faithful-quotient-adjoint-pull}; its restriction at
the zero point of \(G_x\) is zero.
Lemma~\ref{lem:faithful-torus-quotient-seesaw}, with \(Y=G_x\),
gives \(D_a(u,qh)=0\).
Apply that lemma again with \(Y=B\), interchanging the two
variables.  Its pullback is now zero by the last conclusion,
and its zero section is again zero.  Hence \(D_a=0\) on
\(B\times B\).
This proves the actual adjoint line isomorphism for \(M\).
It does not assume that pullback on Picard groups under \(q\)
is injective; the two normalized seesaw applications supply
exactly the needed conclusion.

We describe the primitive normalization in the source
\srcid{ExistsPrimitiveNormalisation}.
Its host element \(w_0\) satisfies
\[
                            w_0+\bar w_0=1.
\]
The theorem of the cube for the line \(M\) on the abelian
variety \(B\) makes \(\lambda_m\) symmetric and biadditive.
Define the line
\[
       P=\mathcal P_M(1,[w_0])
        =[1+[w_0]]^*M\otimes M^{-1}
                         \otimes[w_0]^*M^{-1}\otimes M_0 .
 \tag{AQ3}\label{eq:faithful-primitive-normalized-line}
\]
Thus its pointwise class is \(\lambda_m(z,[w_0]z)\).
The source computes its second difference by the cube identity:
\[
 \begin{split}
 \lambda_{[P]}(u,v)
   &=\lambda_m(u,[w_0]v)+\lambda_m(v,[w_0]u)\\
   &=\lambda_m(u,[w_0]v)+\lambda_m(u,[\bar w_0]v)\\
   &=\lambda_m(u,v).
 \end{split}
 \tag{AQ4}\label{eq:faithful-primitive-polarization}
\]
The middle equality is the adjoint comparison just proved,
together with symmetry.
This is the source's \srcid{diff\_2\_quadratic}.
The same biadditive expansion gives vanishing third difference
for the class in \eqref{eq:faithful-primitive-normalized-line},
and its value at \(-z\) is its value at \(z\).
The defining expression has a canonical identity rigidification.
Choose a generator of the vanishing third-difference line and
normalize it at the identity.
The source's
\srcid{Cube.exists\_normalised}/\srcid{cubicalOfNormalised}
uses global functions of the proper connected group to turn
this normalized generator into the cubical structure:
the two putative cube compatibility expressions differ by a
base unit, whose identity value is \(1\).
The symmetry is the displayed equality at \(-z,z\).
Finally \eqref{eq:faithful-primitive-polarization}, evaluated at
\((z,[w_0]z)\), gives
\[
                              P\simeq\mathcal P_P(1,[w_0]).
\]
This is its primitive isomorphism.
The adjoint isomorphisms for \(P\) follow from
\eqref{eq:faithful-primitive-polarization} and those for \(M\).

The same equality of second differences preserves
\(\mathrm{KTriv}\), as in \srcid{KTrivOfLamEq}.
Indeed, over a local test base, if the translation difference
of \(P\) at \(v\) is trivial, divide out its identity-fibre
constant line to see that \(\lambda_{[P]}(v,-)\) is trivial.
By \eqref{eq:faithful-primitive-polarization},
\(\lambda_m(v,-)\) is trivial.  Restoring the constant fibre
line, which has a generator over the local base, trivializes
the translation difference of \(M\).  Its given
\(\mathrm{KTriv}\) property forces \(v=0\).

It remains essential to recover the original pulled-back line.
The original \(\mathrm{SAB}\) line is already primitive:
\[
                         L_x\simeq\mathcal P_{L_x}(1,[w_0]).
\]
Using \(q^*M\simeq L_x\) and the descended action,
\eqref{eq:faithful-primitive-normalized-line} gives
\[
              q^*P\simeq\mathcal P_{q^*M}(1,[w_0])
                       \simeq\mathcal P_{L_x}(1,[w_0])
                       \simeq L_x .
\]
Thus the normalization on \(B\) produces the required
isomorphism with this same original line.
Together with \(B\)'s properness, the original \(C,i,q\),
and all its kernel and flat-cover fields, this is the
actual \(\mathrm{AbQuot}\) construction.
\end{proof}

\paragraph{Frozen-source ledger and supplied background.}
The assembly, including its two seesaw calls, is
\href{https://github.com/ChristianKudeba/alpoge-rh-verification/blob/alpoge-source-7cbf7a0c/src/Development/AutomorphicSpectrum/Proofs/GeometryProofs/AbelianQuotientOfPowerFns.lean}
{\srcid{AbelianQuotientOfPowerFns}}
(blob \srcid{07fc1095576c3a7ac5bb28ede8f6fea274b8ac1c}):
\srcid{exists\_desc}, \srcid{exists\_act},
\srcid{adj\_quot}, and \srcid{exists\_abQuot}.
The normalized weight/descent proof is
\href{https://github.com/ChristianKudeba/alpoge-rh-verification/blob/alpoge-source-7cbf7a0c/src/Development/AutomorphicSpectrum/Proofs/GeometryProofs/PicardClassTrivialBySeesaw.lean}
{\srcid{PicardClassTrivialBySeesaw}}
(\srcid{fd2dc8406b863f544fe21342a03896050edfd671}),
specifically \srcid{exists\_weight},
\srcid{exists\_unit\_weight}, \srcid{exists\_gen\_of\_weight},
and \srcid{pic\_inj}.
The normalization is
\href{https://github.com/ChristianKudeba/alpoge-rh-verification/blob/alpoge-source-7cbf7a0c/src/Development/AbelianVarietyTheory/Proofs/ModuliProofs/ExistsPrimitiveNormalisation.lean}
{\srcid{ExistsPrimitiveNormalisation}}
(\srcid{16bc270f28fe0061b1a6dadd8c5088713eecb271}),
namely its second-/third-difference identities,
\srcid{exists\_normalised\_bundle}, and final theorem.
These files were checked against the frozen baseline.

The geometric background used here has its exact usual scope:
the theorem of the cube for a line on an abelian variety,
global functions on a proper geometrically connected smooth
group, rigidity of an identity-preserving morphism of abelian
varieties, and fpqc descent of an invertible sheaf with a
generating section.  They are the imported source producers
\srcid{ExistsGenThetaOfIsAbelianScheme},
\srcid{GlobalFnsHolds}, \srcid{AbelianSchemeLawsUnconditional},
and \srcid{CubicalOfNormalised}; the proof above shows where
each is used.  It does not invoke a theorem asserting that the
entire \(\mathrm{AbQuot}\) package exists.
All constructions are on the given geometric fibre and its
actual torus quotient; no integral extension or nonflat
specialization identification is inferred.
